% !TeX program = xelatex
% Fictional fixture: no source-book or course material.
\documentclass[xelatex,ja=standard,11pt,a4paper,openany]{bxjsbook}
\usepackage{amsmath,amsthm}
\usepackage{xcolor}
\usepackage[colorlinks=true,linkcolor=blue]{hyperref}
\usepackage{aliascnt}
\usepackage[nameinlink]{cleveref}
\newtheorem{theorem}{定理}[section]
\numberwithin{equation}{section}
\newaliascnt{lemma}{theorem}
\newtheorem{lemma}[lemma]{補題}
\aliascntresetthe{lemma}
\newaliascnt{proposition}{theorem}
\newtheorem{proposition}[proposition]{命題}
\aliascntresetthe{proposition}
\newaliascnt{corollary}{theorem}
\newtheorem{corollary}[corollary]{系}
\aliascntresetthe{corollary}
\theoremstyle{definition}
\newaliascnt{definition}{theorem}
\newtheorem{definition}[definition]{定義}
\aliascntresetthe{definition}
\newaliascnt{axiom}{theorem}
\newtheorem{axiom}[axiom]{公理}
\aliascntresetthe{axiom}
\newaliascnt{example}{theorem}
\newtheorem{example}[example]{例}
\aliascntresetthe{example}
\newaliascnt{exercise}{theorem}
\newtheorem{exercise}[exercise]{演習}
\aliascntresetthe{exercise}
\crefname{theorem}{\mbox{定理}}{\mbox{定理}}\Crefname{theorem}{\mbox{定理}}{\mbox{定理}}
\crefname{lemma}{\mbox{補題}}{\mbox{補題}}\Crefname{lemma}{\mbox{補題}}{\mbox{補題}}
\crefname{proposition}{\mbox{命題}}{\mbox{命題}}\Crefname{proposition}{\mbox{命題}}{\mbox{命題}}
\crefname{corollary}{\mbox{系}}{\mbox{系}}\Crefname{corollary}{\mbox{系}}{\mbox{系}}
\crefname{definition}{\mbox{定義}}{\mbox{定義}}\Crefname{definition}{\mbox{定義}}{\mbox{定義}}
\crefname{axiom}{\mbox{公理}}{\mbox{公理}}\Crefname{axiom}{\mbox{公理}}{\mbox{公理}}
\crefname{example}{\mbox{例}}{\mbox{例}}\Crefname{example}{\mbox{例}}{\mbox{例}}
\crefname{exercise}{\mbox{演習}}{\mbox{演習}}\Crefname{exercise}{\mbox{演習}}{\mbox{演習}}
\crefname{chapter}{\mbox{章}}{\mbox{章}}\Crefname{chapter}{\mbox{章}}{\mbox{章}}
\crefname{section}{\mbox{節}}{\mbox{節}}\Crefname{section}{\mbox{節}}{\mbox{節}}
\crefname{equation}{\mbox{式}}{\mbox{式}}\Crefname{equation}{\mbox{式}}{\mbox{式}}
\begin{document}
\chapter{架空の参照例}\label{test:chapter}
\section{共有連番}\label{test:section}
\begin{theorem}\label{test:theorem}同じものは同じである。\end{theorem}
\begin{lemma}\label{test:lemma}補題の例。\end{lemma}
\begin{proposition}\label{test:proposition}命題の例。\end{proposition}
\begin{corollary}\label{test:corollary}系の例。\end{corollary}
\begin{definition}\label{test:definition}定義の例。\end{definition}
\begin{axiom}\label{test:axiom}公理の例。\end{axiom}
\begin{example}\label{test:example}具体例。\end{example}
\begin{exercise}\label{test:exercise}演習の例。\end{exercise}
\begin{equation}\label{test:equation}1=1\end{equation}
参照：\cref{test:theorem}、\cref{test:lemma}、\cref{test:proposition}、\cref{test:corollary}、
\cref{test:definition}、\cref{test:axiom}、\cref{test:example}、\cref{test:exercise}。
\Cref{test:lemma}。\cref{test:chapter}、\cref{test:section}、\cref{test:equation}。
\end{document}
